% Solutions to Homework 7 — Math Physics Methods (数学物理方法)

\PassOptionsToPackage{quiet}{xeCJK}
\documentclass[11pt]{article}

\usepackage[margin=1in]{geometry}
\usepackage{amsmath,amsthm,amssymb,mathtools}
\usepackage{ctex}
\usepackage{fancyhdr}
\usepackage{fontspec}
\usepackage{fourier-orns}

\theoremstyle{remark}
\newtheorem{problem}{}
\theoremstyle{definition}
\newtheorem{solution}{解}

\renewcommand{\headrule}{%
\vspace{-8pt}\hrulefill
\raisebox{-2.0pt}{\quad\decofourleft\decotwo\decofourright\quad}\hrulefill}

\setlength{\headheight}{15.11996pt}
\setlength{\topmargin}{-3.11996pt}

\newcommand{\duedate}{\zhdate{2026/11/17}}
\pagestyle{fancy}

\fancyhead{}
\fancyhead[L]{数学物理方法}
\zhnumsetup{style={Traditional,Financial}}
\fancyhead[R]{课后习题柒 \quad 参考解答}

\fancyfoot{}
\fancyfoot[C]{\thepage}
\fancyfoot[L]{截止日期\duedate}
\fancyfoot[R]{任课教师:罗凯}

\begin{document}
\renewcommand{\labelenumi}{(\arabic{enumi})}
\renewcommand{\labelenumii}{(\arabic{enumi}.\arabic{enumii})}

%==========================================================================
\begin{problem}
  将 $\delta(x)$ 展为实数形式的傅里叶积分.
\end{problem}

\begin{solution}
$\delta$ 函数的复数傅里叶积分是
\[
\delta(x)
=\frac{1}{2\pi}\int_{-\infty}^{\infty} e^{\mathrm{i}\omega x}\,\mathrm{d}\omega.
\]
由欧拉公式 $e^{\mathrm{i}\omega x}=\cos(\omega x)+\mathrm{i}\sin(\omega x)$，
\begin{equation*}
\delta(x)
=\frac{1}{2\pi}\int_{-\infty}^{\infty}\cos(\omega x)\,\mathrm{d}\omega
+\frac{\mathrm{i}}{2\pi}\int_{-\infty}^{\infty}\sin(\omega x)\,\mathrm{d}\omega.
\end{equation*}
$\delta(x)$ 是实偶函数，$\sin(\omega x)$ 对 $\omega$ 为奇，故虚部积分为零。于是实数形式为
\[
\delta(x)
=\frac{1}{2\pi}\int_{-\infty}^{\infty}\cos(\omega x)\,\mathrm{d}\omega
=\frac{1}{\pi}\int_0^{\infty}\cos(\omega x)\,\mathrm{d}\omega.
\]
（第二式用了被积函数对 $\omega$ 偶。）
\end{solution}

%==========================================================================
\begin{problem}
  设 $g(x)$ 在实轴上连续，$a\neq 0$ 为常数. 计算
  \[
  \int_{-\infty}^{\infty} g(x)\,\delta\!\left(x^3+ax\right)\,\mathrm{d}x.
  \]
\end{problem}

\begin{solution}
令 $f(x)=x^3+ax=x(x^2+a)$，则 $f'(x)=3x^2+a$。用复合公式
\[
\delta\bigl(f(x)\bigr)
=\sum_{i}\frac{\delta(x-x_i)}{|f'(x_i)|},
\]
其中求和遍及全部简单实零点。

\medskip
\noindent\textbf{情形 $a>0$。}
$x^2+a>0$ 恒成立，唯一实根 $x=0$。且 $f'(0)=a\neq 0$，故
\[
\delta(x^3+ax)=\frac{\delta(x)}{|a|}=\frac{\delta(x)}{a}.
\]
因此
\[
\int_{-\infty}^{\infty} g(x)\,\delta(x^3+ax)\,\mathrm{d}x
=\frac{g(0)}{a}.
\]

\medskip
\noindent\textbf{情形 $a<0$。}
三个简单实根
\[
x_0=0,\qquad
x_{\pm}=\pm\sqrt{-a}.
\]
各点导数值
\begin{align*}
f'(0)&=a,
\qquad |f'(0)|=-a=|a|,\\
f'(x_{\pm})
&=3(-a)+a=-2a=2|a|.
\end{align*}
于是
\[
\delta(x^3+ax)
=\frac{\delta(x)}{|a|}
+\frac{\delta\bigl(x-\sqrt{-a}\bigr)}{2|a|}
+\frac{\delta\bigl(x+\sqrt{-a}\bigr)}{2|a|}.
\]
积分给出
\[
\int_{-\infty}^{\infty} g(x)\,\delta(x^3+ax)\,\mathrm{d}x
=\frac{g(0)}{|a|}
+\frac{g\bigl(\sqrt{-a}\bigr)+g\bigl(-\sqrt{-a}\bigr)}{2|a|}.
\]
\end{solution}

%==========================================================================
\begin{problem}
请阅读 Robin Chapman 写的关于巴塞尔问题的文章《Evaluating $\zeta(2)$》，给出一种其积分表示.
\end{problem}

\begin{solution}
Chapman 文中给出的一种标准积分表示是双重几何级数
\[
\frac{1}{1-xy}
=\sum_{n=0}^{\infty}(xy)^n,
\qquad 0\leq x,y<1.
\]
在单位正方形上积分，
\begin{align*}
\int_0^1\int_0^1\frac{\mathrm{d}x\,\mathrm{d}y}{1-xy}
&=\sum_{n=0}^{\infty}
  \int_0^1\int_0^1 (xy)^n\,\mathrm{d}x\,\mathrm{d}y
=\sum_{n=0}^{\infty}\frac{1}{(n+1)^2}
=\zeta(2).
\end{align*}
因此
\[
\zeta(2)=\int_0^1\int_0^1\frac{\mathrm{d}x\,\mathrm{d}y}{1-xy}.
\]
（与课程中由 $\mathcal{L}\{t\}=1/p^2$ 推出的
$\zeta(2)=\displaystyle\int_0^{\infty}t/(e^t-1)\,\mathrm{d}t$
是同一常数的另一积分写法；二者都等于 $\pi^2/6$。）
\end{solution}

%==========================================================================
\begin{problem}
  计算下列表达式的拉普拉斯变换的像函数
  \begin{enumerate}
    \item $\dfrac{1}{t}\bigl(e^{\beta t}-e^{\alpha t}\bigr)$;
    \item $\displaystyle\int_0^t\dfrac{\sin u}{u}\,\mathrm{d}u$;
    \item $\dfrac{1}{\sqrt{\pi t}}$.
  \end{enumerate}
\end{problem}

\begin{solution}
\noindent\textbf{(1)}
由参数积分规则：若 $f(t)\longleftrightarrow\bar f(p)$，且 $f(t)/t$ 在 $t\to 0$ 附近有界，则
\[
\frac{f(t)}{t}
\longleftrightarrow
\int_p^{\infty}\bar f(q)\,\mathrm{d}q.
\]
取 $f(t)=e^{\beta t}-e^{\alpha t}$，则
$\bar f(p)=1/(p-\beta)-1/(p-\alpha)$（在 $\mathrm{Re}\,p$ 足够大时）。因此
\begin{align*}
\mathcal{L}\Biggl\{\frac{e^{\beta t}-e^{\alpha t}}{t}\Biggr\}
&=\int_p^{\infty}
  \Biggl(\frac{1}{q-\beta}-\frac{1}{q-\alpha}\Biggr)\mathrm{d}q
=\Biggl[\ln\frac{q-\beta}{q-\alpha}\Biggr]_p^{\infty}
=\ln\frac{p-\alpha}{p-\beta}.
\end{align*}
（也可用费曼技巧：令 $I(\lambda)=\int_0^{\infty}e^{-\lambda t}/t\,\mathrm{d}t$ 在差商意义下理解，
对 $\lambda$ 求导得 $-1/\lambda$，再积分并作差。）

\medskip
\noindent\textbf{(2)}
积分定理：$\mathcal{L}\{\int_0^t\psi\}=\mathcal{L}\{\psi\}/p$。取 $\psi(t)=(\sin t)/t$。
由 (1) 的同类计算，或直接
\[
\mathcal{L}\{\sin t\}=\frac{1}{p^2+1}
\qquad\Longrightarrow\qquad
\mathcal{L}\Biggl\{\frac{\sin t}{t}\Biggr\}
=\int_p^{\infty}\frac{\mathrm{d}q}{q^2+1}
=\frac{\pi}{2}-\arctan p.
\]
再除以 $p$，
\[
\mathcal{L}\Biggl\{\int_0^t\frac{\sin u}{u}\,\mathrm{d}u\Biggr\}
=\frac{1}{p}\Biggl(\frac{\pi}{2}-\arctan p\Biggr)
=\frac{\operatorname{arccot} p}{p}.
\]

\medskip
\noindent\textbf{(3)}
先由 $\Gamma$ 函数：
\[
\mathcal{L}\{t^{-1/2}\}
=\int_0^{\infty}t^{-1/2}e^{-pt}\,\mathrm{d}t
=p^{-1/2}\Gamma\Bigl(\tfrac12\Bigr)
=\sqrt{\frac{\pi}{p}}.
\]
因此
\[
\mathcal{L}\Biggl\{\frac{1}{\sqrt{\pi t}}\Biggr\}
=\frac{1}{\sqrt{\pi}}\cdot\sqrt{\frac{\pi}{p}}
=\frac{1}{\sqrt{p}}.
\]
\end{solution}

%==========================================================================
\begin{problem}
  计算下列表达式的拉普拉斯变换的原函数
  \begin{enumerate}
    \item $\bar f(p)=\dfrac{3p}{p^2-1}$;
    \item $\bar f(p)=\dfrac{4p-1}{(p^2+p)(4p^2-1)}$.
  \end{enumerate}
\end{problem}

\begin{solution}
\noindent\textbf{(1)}
裂项
\[
\frac{3p}{p^2-1}
=\frac{3p}{(p-1)(p+1)}
=\frac{3}{2}\Biggl(\frac{1}{p-1}+\frac{1}{p+1}\Biggr).
\]
逐项反演，
\[
f(t)
=\frac{3}{2}\bigl(e^{t}+e^{-t}\bigr)
=3\cosh t.
\]

\medskip
\noindent\textbf{(2)}
分母 $p(p+1)(2p-1)(2p+1)$。设
\[
\frac{4p-1}{p(p+1)(2p-1)(2p+1)}
=\frac{A}{p}+\frac{B}{p+1}+\frac{C}{p-\tfrac12}+\frac{D}{p+\tfrac12}.
\]
留数／待定系数：
\begin{align*}
A
&=\frac{4\cdot 0-1}{(0+1)(0-1)}=\frac{-1}{-1}=1,\\
B
&=\frac{4(-1)-1}{(-1)\bigl(4(-1)^2-1\bigr)}
=\frac{-5}{(-1)\cdot 3}=\frac{5}{3},\\
C
&=\frac{4\cdot\tfrac12-1}{\tfrac12\cdot\tfrac32\cdot 2}
=\frac{1}{\tfrac12\cdot\tfrac32\cdot 2}=\frac{1}{3},\\
D
&=\frac{4\cdot\bigl(-\tfrac12\bigr)-1}{\bigl(-\tfrac12\bigr)\cdot\tfrac12\cdot(-2)}
=\frac{-3}{1}=-3.
\end{align*}
（$C$、$D$ 处把 $(2p\mp 1)$ 的导数因子 $2$ 计入分母。）因此
\[
\bar f(p)
=\frac{1}{p}+\frac{5/3}{p+1}+\frac{1/3}{p-\tfrac12}+\frac{-3}{p+\tfrac12},
\]
原函数
\[
f(t)
=\Theta(t)
+\frac{5}{3}e^{-t}
+\frac{1}{3}e^{t/2}
-3e^{-t/2}.
\]
（$t>0$ 时 $\Theta(t)=1$。）
\end{solution}

%==========================================================================
\begin{problem}
  利用拉普拉斯变换计算
  $\displaystyle\int_0^{\infty}\frac{\sin(xt)}{t}\,\mathrm{d}t$.
\end{problem}

\begin{solution}
对参数 $x$ 讨论。先设 $x>0$。利用
\[
\frac{1}{t}=\int_0^{\infty}e^{-st}\,\mathrm{d}s
\qquad(t>0),
\]
交换积分次序（$t>0$ 时合法），
\begin{align*}
\int_0^{\infty}\frac{\sin(xt)}{t}\,\mathrm{d}t
&=\int_0^{\infty}\sin(xt)
  \Biggl(\int_0^{\infty}e^{-st}\,\mathrm{d}s\Biggr)\mathrm{d}t
=\int_0^{\infty}\mathrm{d}s
  \int_0^{\infty}\sin(xt)\,e^{-st}\,\mathrm{d}t
\\
&=\int_0^{\infty}\frac{x}{s^2+x^2}\,\mathrm{d}s
=\Biggl[\arctan\frac{s}{x}\Biggr]_0^{\infty}
=\frac{\pi}{2}.
\end{align*}
当 $x=0$ 时被积函数恒为零，积分为 $0$。
当 $x<0$ 时令 $x=-|x|$，
\[
\int_0^{\infty}\frac{\sin(xt)}{t}\,\mathrm{d}t
=-\int_0^{\infty}\frac{\sin(|x|t)}{t}\,\mathrm{d}t
=-\frac{\pi}{2}.
\]
合并即
\[
\int_0^{\infty}\frac{\sin(xt)}{t}\,\mathrm{d}t
=\frac{\pi}{2}\operatorname{sgn}(x)
=
\begin{cases}
\pi/2, & x>0,\\
0, & x=0,\\
-\pi/2, & x<0.
\end{cases}
\]
\end{solution}

%==========================================================================
\begin{problem}
用拉普拉斯反变换中的 Bromwich（又称黎曼--梅林）积分求
\[
\frac{p}{p^2-\omega^2}
\]
的原函数.
\end{problem}

\begin{solution}
设 $\omega$ 为实常数，$\omega\neq 0$。Bromwich 积分
\[
f(t)
=\frac{1}{2\pi\mathrm{i}}
\int_{\sigma-\mathrm{i}\infty}^{\sigma+\mathrm{i}\infty}
\frac{p\,e^{pt}}{p^2-\omega^2}\,\mathrm{d}p,
\qquad t>0,
\]
其中 $\sigma>|\omega|$，使竖直线位于全部极点右侧。被积函数在
$p=\pm\omega$ 有简单极点。对 $t>0$ 向左半平面闭上（约当引理），
\begin{align*}
f(t)
&=\mathrm{Res}_{p=\omega}
  \Biggl[\frac{p\,e^{pt}}{(p-\omega)(p+\omega)}\Biggr]
+\mathrm{Res}_{p=-\omega}
  \Biggl[\frac{p\,e^{pt}}{(p-\omega)(p+\omega)}\Biggr]
\\
&=\frac{\omega\,e^{\omega t}}{2\omega}
+\frac{(-\omega)\,e^{-\omega t}}{-2\omega}
=\frac{e^{\omega t}+e^{-\omega t}}{2}
=\cosh(\omega t).
\end{align*}
与查表 $\mathcal{L}\{\cosh(\omega t)\}=p/(p^2-\omega^2)$ 一致。
\end{solution}

%==========================================================================
\begin{problem}
  利用拉普拉斯变换求解微分方程
  \[
  \frac{\mathrm{d}^3 y}{\mathrm{d}t^3}
  +3\frac{\mathrm{d}^2 y}{\mathrm{d}t^2}
  +3\frac{\mathrm{d} y}{\mathrm{d}t}
  +y
  =6e^{-t},
  \qquad
  y(0)=y'(0)=y''(0)=0.
  \]
\end{problem}

\begin{solution}
记 $Y(p)=\mathcal{L}\{y\}$。初值全为零，故
\[
\mathcal{L}\{y^{(n)}\}=p^n Y(p),\qquad n=1,2,3.
\]
右边 $\mathcal{L}\{6e^{-t}\}=6/(p+1)$。变换后的方程是
\[
\bigl(p^3+3p^2+3p+1\bigr)Y(p)=\frac{6}{p+1}.
\]
左边括号正是 $(p+1)^3$，因此
\[
Y(p)
=\frac{6}{(p+1)\,(p+1)^3}
=\frac{6}{(p+1)^4}.
\]
由公式
\[
\mathcal{L}\bigl\{t^n e^{at}\bigr\}
=\frac{n!}{(p-a)^{n+1}}
\qquad\Longrightarrow\qquad
\mathcal{L}^{-1}\Biggl\{\frac{1}{(p-a)^{n+1}}\Biggr\}
=\frac{t^n e^{at}}{n!},
\]
取 $a=-1$、$n=3$，
\[
\mathcal{L}^{-1}\Biggl\{\frac{1}{(p+1)^4}\Biggr\}
=\frac{t^3 e^{-t}}{3!}=\frac{t^3 e^{-t}}{6}.
\]
故
\[
y(t)=6\cdot\frac{t^3 e^{-t}}{6}=t^3 e^{-t}.
\]
\end{solution}

\end{document}
